Fix \(k=(g,t)\in\mathcal K\) and \(c\in\mathcal J_k\). For \(a\in\{g,c\}\) and \(s\in\{0,1\}\), define the untreated conditional change
\[
 \mu^0_{a,s,k}(x)
 =E_Q\!\left[Y_{it}(\infty,s)-Y_{i,g-1}(\infty,s)
       \mid G_i=a,S_i=s,X_i=x\right].
\]
Also define the target-cell conditional effect
\[
 q_k(x)=E_Q\!\left[Y_{it}(g,1)-Y_{it}(\infty,1)
       \mid G_i=g,S_i=1,X_i=x\right].
\]
Consistency and no anticipation at the observed baseline \(g-1<g\) give
\[
 m_{g,1,k}(x)=q_k(x)+\mu^0_{g,1,k}(x).
\]
No spillover to ineligible units gives
\[
 m_{g,0,k}(x)=\mu^0_{g,0,k}(x).
\]
Because \(c>t\ge g\), both \(g-1<c\) and \(t<c\). Consistency and no anticipation therefore give
\[
 m_{c,1,k}(x)=\mu^0_{c,1,k}(x),
 \qquad
 m_{c,0,k}(x)=\mu^0_{c,0,k}(x),
\]
where the second equality also follows from no spillover. Substitution into the observed conditional DDD contrast yields
\[
 d_{k,c}(x)=q_k(x)
 +\{\mu^0_{g,1,k}(x)-\mu^0_{g,0,k}(x)\}
 -\{\mu^0_{c,1,k}(x)-\mu^0_{c,0,k}(x)\}.
\]
The pairwise conditional DDD restriction makes the last two braces equal for \(Q_X\)-almost every \(x\). Hence
\[
 d_{k,c}(x)=q_k(x)
\]
almost surely on the covariate support relevant to the target cell. By iterated expectation and the definition of the generalized propensity,
\[
\begin{aligned}
 \frac{E_P[p_{g,1}(X_i)d_{k,c}(X_i)]}{\pi_{g,1}}
 &=\frac{E_P[\mathbf 1\{G_i=g,S_i=1\}q_k(X_i)]}{\pi_{g,1}}\\
 &=E_Q[Y_{it}(g,1)-Y_{it}(\infty,1)\mid G_i=g,S_i=1]\\
 &=\theta_k(Q).
\end{aligned}
\]
Overlap makes \(\pi_{g,1}=E_P[p_{g,1}(X_i)]\ge\varepsilon>0\), so every division above is valid. The equality holds for every \(c\in\mathcal J_k\); averaging over \(c\) therefore gives \(\Theta_k(P)=\theta_k(Q)\).

It remains to verify the score mean. Conditional on \((G_i,S_i,X_i)\), every residual \(\Delta Y_{i,k}-m_{a,s,k}(X_i)\) appearing in the score has mean zero. Thus all residual summands in \(R_{k,c}\) integrate to zero. Moreover,
\[
 E_P\!\left[\frac{\mathbf 1\{G_i=g,S_i=1\}}{\pi_{g,1}}
 \{d_{k,c}(X_i)-\Theta_k(P)\}\right]
 =\frac{E_P[p_{g,1}(X_i)d_{k,c}(X_i)]}{\pi_{g,1}}-\Theta_k(P)=0.
\]
Consequently \(E_P[R_{k,c}(O_i;\eta_0)]=\Theta_k(P)=\theta_k(Q)\). Stacking the coordinate equalities gives
\[
 \theta(Q)=\Theta(P),\qquad
 E_P[R(O_i;\eta_0)]=B\Theta(P),\qquad
 E_P[\Psi(O_i)]=0.
\]
Premultiplication by the prespecified \(\omega\) gives
\[
 \tau_\omega(Q)=\omega^{\prime}\theta(Q)
 =\omega^{\prime}\Theta(P)=\Theta_\omega^{\mathrm{agg}}(P).
\]
Finally, if \(g=2\), then \(g-1=1\), so every change used above is the observed change from period \(1\) to period \(t\).